% 张三慧《大学物理学》第4版 V2，第22章《光的干涉》。
% 主教材PDF物理页262--264（印刷245--247），解答S PDF312--318。
% 原书图22.33--22.37待矢量重绘与专家图示审查。
\section*{第22章\quad 光的干涉}
\subsection*{章末习题22.1--22.23}
\paragraph{习题 22.1}
钠黄光 \(\lambda=589.3\,\mathrm{nm}\)，一次发光持续 \(10^{-8}\,\mathrm s\)。求一个波列包含的振动周期数。
\paragraph{习题 22.2}
汞灯发光经过滤光片照射双缝，缝距 \(d=0.60\,\mathrm{mm}\)，屏距 \(D=2.5\,\mathrm m\)，相邻明纹距 \(\Delta x=2.27\,\mathrm{mm}\)。求波长和颜色。
\paragraph{习题 22.3}
劳埃德镜（图22.33）：光源高出镜面 \(h=2\,\mathrm{mm}\)，源至镜左端水平间距 \(20\,\mathrm{cm}\)，镜长 \(30\,\mathrm{cm}\)，屏在镜右端。入射波长 \(\lambda=7.2\times10^{-7}\,\mathrm m\)，求镜右端与屏上第一明纹的距离。需考虑镜面反射相位跃变。
\paragraph{习题 22.4}
双缝间距 \(0.15\,\mathrm{mm}\)，屏距 \(1.0\,\mathrm m\)，第1级到第10级明纹距离为 \(36\,\mathrm{mm}\)。求波长。
\paragraph{习题 22.5}
南北方向相距 \(3.0\,\mathrm{km}\) 的两无线电台同时发射频率 \(2.0\times10^5\,\mathrm{Hz}\) 的电波，南台比北台相位落后 \(\pi/2\)。求远场相长方向的方位角（相对东西方向，北偏为正）。
\paragraph{习题 22.6}
两相干水波源间距 \(12\,\mathrm{cm}\)，在前方 \(50\,\mathrm{cm}\) 处相邻平静区中心相距 \(4.5\,\mathrm{cm}\)，波速 \(25\,\mathrm{cm/s}\)。求频率。
\paragraph{习题 22.7}
光以入射角 \(\varphi\) 倾斜射向双缝，缝距 \(d\)。(1) 证明明纹出射角满足 \(d\sin\theta-d\sin\varphi=\pm k\lambda\)，\(k=0,1,\ldots\)。(2) 证明小出射角时相邻明纹角间隔 \(\Delta\theta\) 与 \(\varphi\) 无关。
\paragraph{习题 22.8}
澳大利亚早期天文射电观测：一束 \(f=6.0\times10^7\,\mathrm{Hz}\) 电波直接到天线，另一束经海面反射到天线（图22.34），天线高 \(25\,\mathrm m\)。求相消干涉对应海面\emph{最小非零掠射角} \(\theta\)。
\paragraph{习题 22.9}
证明等强双缝干涉明纹按半峰值定义的角宽度 \(\Delta\theta=\lambda/(2d)\)，其中宽度指明纹两侧强度为中心强度一半的位置之间的全宽。
\paragraph{习题 22.10}
图22.35：用激光与分光镜 \(M_1\)、反射镜 \(M_2\) 形成两重叠光束，光束与屏夹角为 \(45^\circ\)，振幅比 \(A_1:A_2=2:1\)，波长 \(632.8\,\mathrm{nm}\)。求屏上干涉纹间距和可见度。
\paragraph{习题 22.11}
氦氖激光器发射 \(\lambda=632.8\,\mathrm{nm}\) 红光，谱线频率宽度 \(\Delta\nu=1.3\times10^9\,\mathrm{Hz}\)。求相干长度与相干时间。
\paragraph{习题 22.12}
太阳地面视角宽度约 \(10^{-2}\,\mathrm{rad}\)，光波长 \(550\,\mathrm{nm}\)。用太阳光进行双缝干涉，估算缝距不能超过的空间相干间隔。
\paragraph{习题 22.13}
双缝的一条缝盖薄玻璃，原第7级明纹移到原零级明纹位置。 \(\lambda=550\,\mathrm{nm}\)，玻璃折射率 \(n=1.58\)。求玻璃厚度。
\paragraph{习题 22.14}
Si衬底折射率 \(3.42\)、SiO\(_2\) 薄膜折射率 \(1.50\)，薄膜蚀出薄楔（图22.36），正入射 \(\lambda=589.3\,\mathrm{nm}\) 观察到从零厚处起第7条反射暗纹。求此处SiO\(_2\)膜厚 \(h\)，并说明反射相位损失。
\paragraph{习题 22.15}
厚 \(0.4\,\mathrm{\mu m}\)、折射率 \(1.50\) 的玻璃薄片，白光正入射。求可见光范围内反射加强、透射加强的波长。取可见范围约400--760\,nm。
\paragraph{习题 22.16}
为了使折射率 \(1.5\) 的人造水晶在波长 \(560\,\mathrm{nm}\) 有强反射，在其表面镀折射率 \(2.0\) 的一氧化硅薄膜。求至少镀多厚。
\paragraph{习题 22.17}
\(n=1.5\) 的玻璃表面覆盖折射率 \(1.32\) 的油膜，正入射连续可调单色光。波长485nm时反射相消；波长增至679nm时\emph{再次}相消。求油膜厚度（两次消光为相邻级次）。
\paragraph{习题 22.18}
白光照射折射率 \(1.33\) 的肥皂膜，从与法线成 \(45^\circ\) 的方向看反射呈绿色 \(\lambda=500\,\mathrm{nm}\)。求最小膜厚；从法线方向看呈什么颜色？假设薄膜两侧为空气。
\paragraph{习题 22.19}
折射率 \(n_1=1.52\) 的镜头表面镀 \(n_2=1.38\) 的 MgF\(_2\) 增透膜，适用波长 \(550\,\mathrm{nm}\)。求最小厚度。
\paragraph{习题 22.20}
单色光牛顿环，一明环直径 \(3.00\,\mathrm{mm}\)，其外面第5个明环直径 \(4.60\,\mathrm{mm}\)，平凸透镜曲率半径 \(R=1.03\,\mathrm m\)。求光波长。
\paragraph{习题 22.21}
折射率 \(n\)、厚度 \(h\) 的薄玻璃片放在迈克耳孙干涉仪一臂，问两光路光程差改变多少。
\paragraph{习题 22.22}
迈克耳孙干涉仪可动镜移动 \(\Delta L=0.3220\,\mathrm{mm}\)，等倾干涉纹在中心缩进1204条。求波长。
\paragraph{习题 22.23}
气体折射率干涉仪（图22.37）两臂各有长 \(d=20\,\mathrm{cm}\) 的平行玻璃气管，先抽真空，后在其中一管逐渐充入气体到标准状态；波长 \(589.3\,\mathrm{nm}\)，经过98条干涉纹。求气体折射率。需按原图确认每束通过气管的次数。
